% Stand 26. Juli 2026

\subsection{Verwandte Arbeiten}
\label{KapitelDatenauswertung_VerwandteArbeiten}
Es existieren themenverwandte Arbeiten, die überwiegend aber
eine andere Zielsetzung verfolgen:
\begin{itemize}[topsep=2pt,parsep=1pt,itemsep=0.5pt]
\item[$\circ$] In der Arbeit \cite{Kim2014} werden
die jährlichen Mittelwerte der 
$\mathrm{PM}_{10}$-Konzen\-tra\-tion in koreanischen 
Großstädten mittels ordinary Kriging untersucht. 
Insgesamt werden 226 Messstationen untersucht. Die 
durch Schätzung der empirischen Semivarianz ermittelten
Werte für den Range schwanken zwischen  $8.8\,km$ für Seoul
und $29.5\,km$ für das gesamte untersuchte Areal.
Die Autoren konstatieren schlechte Resultate basierend auf 
Methode und Datenlage.
\item[$\circ$] In der Arbeit \cite{Jang2023} werden
die jährlichen Mittelwerte der 
$\mathrm{PM}_{2.5}$-Konzentration in Taiwan mithilfe des
in dieser Arbeit nicht beschriebenen multivariaten 
Indikator Krigings untersucht. 
In jener Arbeit werden für die Jahre 2019-2021 Ranges
zwischen $30-50\,km$ und für die Jahre 2020-2022 
Ranges von $20-60\,km$ ermittelt.\\
Die von den Autoren ermittelte Spannbreite für den Range unterstützt die 
Feststellung, dass der Range im Rahmen dieser Arbeit prinzipiell gar 
nicht schätzbar wäre, weil die maximale paarweise Distanz der 
untersuchten Sensoren zu gering ist.
\item[$\circ$] In der Arbeit \cite{Shukla2020} werden 
die monatlichen und jährlichen Mittelwerte der
$\mathrm{PM}_{2.5}$-Konzentration von rund 30 Messstationen
innerhalb der indischen Stadt Delhi mithilfe des ordinary
Kriging und der in dieser Arbeit nicht beschriebenen 
Methode des Inverse Distance Weighting vergleichend 
untersucht. In jener Studie liefern beide Methoden 
vergleichbare Ergebnisse. Die Autoren merken an, dass
die für die Anwendung des Kriging wichtige 
Voraussetzung der Stationarität möglicherweise 
nicht erfüllt sein könnte.  
\item[$\circ$] In der Arbeit \cite{Ramos2016} werden die
täglichen Mittelwerte der $\mathrm{PM}_{2.5}$-Konzen\-tra\-tion
auf der kanadischen Insel Montreal mithilfe 
der Methoden des Inverse Distance Weighting und 
des Kriging mit externem Drift untersucht. Der Fokus
jener Arbeit liegt auf einer möglichst guten 
raum-zeitlichen Vorhersage der 
$\mathrm{PM}_{2.5}$-Konzentra\-ti\-onen und die Autoren
favorisieren hierfür ein hybrides Modell aus 
Inverse Distance Weighting und Kriging mit Drift unter Einbeziehung von Kovariablen. 
Für jene Untersuchung werden die Messungen von rund 
10 Stationen herangezogen.
\end{itemize}
